\documentclass[a4paper,UKenglish,cleveref,autoref]{lipics-v2021}

\hideLIPIcs

\bibliographystyle{plainurl}% the mandatory bibstyle

\title{A Simulated Annealing Heuristic for Maximum Agreement Forest}

\author{Adam Polak}{Bocconi University}{adam.polak@unibocconi.it}{}{}
\authorrunning{A. Polak}
\Copyright{Adam Polak}

\ccsdesc[500]{Theory of computation~Design and analysis of algorithms}
\ccsdesc[300]{Applied computing~Bioinformatics}

\keywords{maximum agreement forest, rooted phylogenetic trees, simulated annealing, PACE}

\supplement{\url{https://github.com/adampolak/samaf}}

\begin{document}
\maketitle

\begin{abstract}
We describe a solver for the Maximum Agreement Forest submitted for
the heuristic track of PACE 2026 Challenge.
The solver uses simulated annealing search representing a candidate solution by
the set of edges cut in one input tree.
Moreover it employs a few known reduction rules from the literature,
and some ad-hoc pre- and post-processing heuristics.
\end{abstract}

\section{Overview}

The input consists of two rooted binary phylogenetic trees \(T_1,T_2\) on the
same taxon set \(\{1,\ldots,n\}\), given in the PACE Newick format.  The output
is a set of rooted trees, one per line, that is obtainable from both inputs by
deleting directed edges and suppressing degree-one vertices.  The objective is
to minimize the number of output components.  For two rooted trees this is the
usual rooted maximum agreement forest problem, equivalently the rooted SPR
distance plus one under the PACE convention \cite{pace2026maf,bordewich2005}.

The solver stores each parsed tree in postorder: every child has smaller index
than its parent.  This makes bottom-up dynamic programs a plain forward scan and
top-down propagation a reverse scan.  Preprocessing stores parent and child
indices, leaf-to-node maps, Euler intervals, leaf ranks, subtree sizes, binary
lifting ancestors, and a canonical identifier for every rooted subtree.  The
canonical identifiers are maintained by an interner over ordered pairs of child
identifiers.  They let the implementation compare cleaned rooted shapes by
integer equality instead of string comparison.

\section{Cut Encoding and Evaluation}

A search state is a bit mask \(c\) on the non-root edges of \(T_1\).  Edge
\(e\) is cut iff \(c_e=1\); the connected components of the remaining forest
induce the candidate partition.

Given \(c\), the evaluator first assigns every node of \(T_1\) to the
root of its cut-induced component.  It then computes, for each component, the
canonical cleaned shape in \(T_1\) and its total number of leaves.  A second
linear scan over \(T_2\) checks whether each component appears as one connected
cleaned subtree and whether different components ever overlap on an edge.  At
any \(T_2\) node there may be at most one unfinished component continuing
upward; two different unfinished components at the same node give an immediate
conflict.  A component is feasible iff the completed shape in \(T_2\) matches
the shape computed from \(T_1\).  The evaluator returns the component count and
a penalty count.  The penalty count is zero exactly for valid agreement
forests.  Positive penalty is accumulated for local certificates of
infeasibility: a leaf assigned to no component, two different open components
meeting at the same \(T_2\) edge, an unfinished component reaching the root of
\(T_2\), a component that never closes in \(T_2\), or a closed component whose
canonical \(T_2\) shape differs from its \(T_1\) shape.  The annealing energy is
\[
  E(c)=\#\mathrm{components}(c)+8\cdot \#\mathrm{penalties}(c).
\]
The constant $8$ in the above expression, as many other constants in the solver,
were chosen empirically.

The hot loop uses a fast evaluator with the same control flow, but replaces
canonical shape identifiers by 64-bit polynomial hashes.  This has the same
linear-time asymptotic complexity but in practice this is much cheaper
than interning every temporary shape during annealing.  To avoid returning a
hash false positive, every incumbent improvement is rechecked by the exact
evaluator, and the final components are emitted only from an exact feasible
evaluation.

\section{Preprocessing: Reductions and Decomposition}

Before search, the solver applies exhaustive local reductions.
\emph{Common cherries} repeatedly contract identical pendant subtrees
\cite{bordewich2005}.
\emph{Oriented common four-chain} reductions delete the fourth taxon 
of a matching rooted chain and remember how to expand it
\cite{bordewich2005,KelkLM23}.
The implemented \emph{3--2 chain} reduction removes a forced singleton component;
this is the parameter-changing rule from the rooted kernelization literature
\cite{KelkLM23}.  The reduced instance is relabelled densely, solved, and then
expanded back to the original labels.  Expansion first lifts reduced components
through the representative map, then reinserts deleted chain taxa into the
appropriate component when possible, and finally appends forced singleton
components.

Large instances are recursively split on \emph{common clusters}
\cite{bordewich2005,linzsemple2011,LiZ17}.  A cluster is found as a
\(T_1\) subtree whose leaf ranks form an interval in \(T_2\) and whose size is
large enough on both sides.  The most balanced such split is preferred.  The
exact common-cluster reduction has an off-by-one case: depending on whether an
optimal forest contains a component bridging the cluster boundary, the original
optimum is either the sum of the two subproblem optima or one less.  Exact
implementations resolve this by solving additional variants of the split
subinstances.  Our heuristic does not solve these variants.  It simply solves
the inside and outside subinstances independently, accepting the possible loss
of one component per split in exchange for much smaller search spaces.  A final
global merge reparses the original instance and tries to recover such
boundary-bridging components by merging components across the split.

There is also an early cut heuristic for very large blocks.  Before committing
to a binary cluster split, the solver may run a short annealing phase on the
block.  If this already finds a feasible forest with many components, the code
searches for up to two disjoint common clusters whose sizes are close to
\(n/3\), with a minimum side-size constraint.  These clusters and the remaining
outside part become several children at once.  This is not a reduction rule; it
is a decomposition heuristic that spends a small amount of search time to find
a more balanced set of subinstances.

\section{Initialization}

The initial feasible solution is set to the singleton forest.
On small instances ($n \leq 2000$) the solver also spends a small fraction of
the budget on trying to greedily improve this initial solution before the
annealing phase. The greedy improvement consists of going through all the cut
edges one by one and uncutting them if this does not make the solution infeasible.

\section{Annealing}

The annealing temperature follows an exponential schedule from \(0.60\) to \(0.012\),
following the standard simulated annealing paradigm \cite{kirkpatrick1983}.
A proposed move is applied in-place and recorded as a short rollback list of changed edges.
If the Metropolis test rejects the move, the rollback list restores the previous cut
mask.

The move distribution is biased toward reducing the number of cut
edges: single toggle, remove, add, pair-remove, interval-guided swap, sibling
swap, and parent/child level swap.
The toggle move flips one random non-root edge.  The remove move samples cut
edges and unsets one of them; add does the symmetric operation on an uncut edge.
Pair-remove first removes one cut edge and then tries to remove a second cut
edge chosen near it in a precomputed leaf-rank order.  The interval-guided swap
removes a cut edge and tries to add an uncut edge with nearby \(T_2\) ranks,
falling back to a plain removal if no such edge is found.  The sibling swap
moves a cut from one child edge of a \(T_1\) fork to the sibling edge.  The
level swap moves a cut one level upward to the parent edge, or one level
downward to one of its child edges.

Swap moves are more important at low temperature, so their weight is increased
near the end of the schedule.
Several proposal types use precomputed edge orders by the minimum and maximum leaf rank
of a \(T_1\) subtree in both \(T_1\) and \(T_2\).  This gives local moves that
tend to exchange cuts between nearby intervals rather than random unrelated
edges.

Because reductions, cluster decomposition, and the early cut heuristic can
produce several independent leaf blocks, the solver does not finish one block
before starting the next.  It initializes every leaf block, reserves time for
polishing and the final global merge, and then interleaves annealing steps in
rounds.  In each step a block is sampled with weight proportional to its current
component count.  This is a sequential simulation of parallel search: every block
keeps an exact feasible incumbent, so interruption leaves a complete solution.

\section{Postprocessing}

After annealing, the solver spends the reserved polishing budget ($1$ second) on
deterministic improvements.
Greedy deletion tries to remove cut edges while preserving feasibility.
A swap polish tries sibling, ancestor, and descendant cut exchanges and then
repeats deletion.

\section{Development}

This solver was developed as an experiment with heavy use of coding agents,
primarily OpenAI Codex with GPT-5.5 and Claude Code with Sonnet 4.6.
They were used for literature review, implementing initial versions of
multiple approaches, and engineering improvements and heuristics in a continuous
loop.

We estimate that this project took 40 hours of human time, while the coding agents
run for 140 hours in total and were given access to 50 CPU cores for testing.
The API calls to the language models used approximately 250 million tokens.

\bibliography{references}

\end{document}
